% PGelisa på tur – fargeprofil for PDF (xelatex eller lualatex, KOMA-Script).
% Fargene er rollene i designsystemet «PGelisa på tur».

\usepackage{xcolor}
\usepackage{graphicx}

% Flater og tekst
\definecolor{pgskum}{HTML}{F6F9F7}
\definecolor{pgflate}{HTML}{FFFFFF}
\definecolor{pgflatetonet}{HTML}{DCF0EC}
\definecolor{pgtekst}{HTML}{0B2B2E}
\definecolor{pgtekstdempet}{HTML}{4F5E5D}
\definecolor{pglinje}{HTML}{D2D9D8}
\definecolor{pglinjesterk}{HTML}{6B7B79}
% Merkefarger
\definecolor{pgmerke}{HTML}{13A091}
\definecolor{pgmerketekst}{HTML}{06695F}
\definecolor{pgdyp}{HTML}{0E4A52}
\definecolor{pgaksentlav}{HTML}{C9D43C}
% Status
\definecolor{pgok}{HTML}{06695F}
\definecolor{pgokflate}{HTML}{EDFAF8}
\definecolor{pginfo}{HTML}{185E95}
\definecolor{pginfoflate}{HTML}{F1F8FE}
\definecolor{pgadvarsel}{HTML}{9A3501}
\definecolor{pgadvarselflate}{HTML}{FFF4F0}
\definecolor{pgfeil}{HTML}{B81D2E}
\definecolor{pgfeilflate}{HTML}{FEF4F3}

% Fonter: Atkinson Hyperlegible i brødtekst, Bricolage Grotesque i titler.
% Stiene gjelder dokumenter som ligger i prosjektets rotmappe.
\usepackage{fontspec}
\setmainfont{AtkinsonHyperlegible}[
  Path = _extensions/pgelisa/fonts/,
  Extension = .ttf,
  UprightFont = *-Regular,
  BoldFont = *-Bold,
  ItalicFont = *-Italic,
  BoldItalicFont = *-BoldItalic]
\setsansfont{AtkinsonHyperlegible}[
  Path = _extensions/pgelisa/fonts/,
  Extension = .ttf,
  UprightFont = *-Regular,
  BoldFont = *-Bold,
  ItalicFont = *-Italic,
  BoldItalicFont = *-BoldItalic]
\newfontfamily\pgtitlerfont{BricolageGrotesque}[
  Path = _extensions/pgelisa/fonts/,
  Extension = .ttf,
  UprightFont = *-Bold,
  BoldFont = *-Bold]
\newfontfamily\pgtitlerfonthalv{BricolageGrotesque}[
  Path = _extensions/pgelisa/fonts/,
  Extension = .ttf,
  UprightFont = *-SemiBold,
  BoldFont = *-Bold]

\color{pgtekst}

% Tegn som mangler i Atkinson Hyperlegible
\usepackage{newunicodechar}
\newunicodechar{→}{\ensuremath{\rightarrow}}
\newunicodechar{←}{\ensuremath{\leftarrow}}

% Lange kodelinjer brytes i stedet for å gå ut over kanten
\usepackage{fvextra}
\AtBeginDocument{%
  \RecustomVerbatimEnvironment{Highlighting}{Verbatim}{breaklines,breakanywhere,commandchars=\\\{\}}%
}

% Callout-farger i profilens statusfarger. Quarto definerer sine egne etter denne fila,
% så de overstyres når dokumentet starter.
\AtBeginDocument{%
  \definecolor{quarto-callout-note-color}{HTML}{185E95}%
  \definecolor{quarto-callout-note-color-frame}{HTML}{185E95}%
  \definecolor{quarto-callout-tip-color}{HTML}{06695F}%
  \definecolor{quarto-callout-tip-color-frame}{HTML}{06695F}%
  \definecolor{quarto-callout-important-color}{HTML}{0B2B2E}%
  \definecolor{quarto-callout-important-color-frame}{HTML}{0B2B2E}%
  \definecolor{quarto-callout-warning-color}{HTML}{9A3501}%
  \definecolor{quarto-callout-warning-color-frame}{HTML}{9A3501}%
  \definecolor{quarto-callout-caution-color}{HTML}{B81D2E}%
  \definecolor{quarto-callout-caution-color-frame}{HTML}{B81D2E}%
}

% Overskrifter og tittel (KOMA-Script)
\makeatletter
\@ifclassloaded{scrartcl}{%
  \setkomafont{disposition}{\pgtitlerfont\color{pgtekst}}
  \setkomafont{title}{\pgtitlerfont\color{pgdyp}}
  \setkomafont{subtitle}{\pgtitlerfonthalv\color{pgtekstdempet}\large}
  \setkomafont{author}{\normalfont\color{pgtekstdempet}}
  \setkomafont{date}{\normalfont\color{pgtekstdempet}}
  \titlehead{\includegraphics[height=1.3cm]{_extensions/pgelisa/logo/pgelisa-logo-liggende.pdf}}
}{}
\makeatother

% Taster (kbd-snarveien) i en tynn ramme
\usepackage[skins,breakable]{tcolorbox}
\newtcbox{\pgkbd}{on line, colback=pgflate, colframe=pglinjesterk,
  boxrule=0.4pt, arc=2pt, left=2pt, right=2pt, top=0.5pt, bottom=0.5pt, boxsep=0pt}
\providecommand{\kbd}[1]{\pgkbd{\small #1}}
